Los entrenamientos fueron agrupados por funciones de pérdida, en particular por una estructura de tipos dentro de estas para ayudar a su posterior documentación y evaluación:

\begin{itemize}
    \item \textbf{Pérdidas principales}: $L1$, $L2$, $L1\_freq$, $\log L1$, $\log L2$, $\log\_mag$, $\log\_compressed\_l2$, $LPSA$, $mask\_L1$.
    \item \textbf{Pérdidas experimentales}: deep\_feature, deep\_feature\_emd.
    \item \textbf{Pérdidas avanzadas}: l\_mrs, lpsa\_phase.
\end{itemize}

También existe la división que ya se explicó anteriormente entre dos y cuatro fuentes. El caso de las dos fuentes representa una separación más sencilla y cercana a escenarios de voz y acomppañamiento.
El caso de cuatro fuentes introduce mayor compplpepjidad, ya que el modelo debe repppartir la mezcla entre varias fuentes instrumentales con solapamiento espectral.
Se respetó la estructura que presenta el \textit{dataset} utilizado de \textit{MUSDB18}, compuesto por 150 pistas con \textit{stems} aislados, con etiquetas en el orden de \textit{vocales,baterías,bajos y otros} y dividido en conjuntos de entrenamiento y test.